% Replacement for Sections 5.1--5.3 of Emergence Dynamics v0.2.
% This text proves a representation/embedding result. It does not claim
% computational power beyond the Church--Turing framework.

\subsection{The Measure-Dynamical Envelope and Its Non-Degenerate Sector}

To state precisely in what sense a classical Turing machine is a degenerate
case, we must distinguish the full measure-dynamical framework from its
strictly emergent sector.  Strict uncertainty, state-space expansion, and
metric deformation cannot be axioms imposed on every member of the framework:
a deterministic computation has zero state uncertainty and a fixed transition
law.  We therefore use the following inclusive definition.

\begin{definition}[Measure-dynamical computation]
 A measure-dynamical computation is a tuple
 \[
   \mathfrak E=(X,\mathcal B,\{K_t\}_{t\in T},\{\mu_t\}_{t\in T}),
 \]
 where $(X,\mathcal B)$ is a standard Borel state space,
 $\mu_t\in\mathcal P(X)$ is the computational state law, and $K_t$ is a
 Markov transition kernel.  In discrete time,
 \[
   \mu_{k+1}(A)=\int_X K_k(x,A)\,d\mu_k(x),
   \qquad A\in\mathcal B.
 \]
 The framework may additionally carry an information filtration
 $\{\mathcal I_t\}$, a time-dependent metric $g_t$, and endogenous parameters
 $\Theta_t$ governing the kernels.
\end{definition}

\begin{definition}[Non-degenerate emergence and the classical boundary]
 A measure-dynamical computation is called non-degenerately emergent on an
 interval if at least one of the following occurs on that interval:
 \begin{enumerate}
   \item $K_t(x,\cdot)$ is non-Dirac on a set of positive $\mu_t$-measure;
   \item a diffusion or jump-dispersion term is non-zero;
   \item the transition law, metric, or effective state representation changes
         endogenously with incoming information;
   \item the state law has positive entropy in a specified discrete
         coarse-graining or develops non-atomic support.
 \end{enumerate}
 The classical deterministic boundary is obtained when all four mechanisms
 vanish: the kernel is time-independent and Dirac-valued, diffusion is zero,
 the state representation is fixed, and the state law remains atomic.
\end{definition}

Accordingly, the strict statements $H(\mu_t\mid\mathcal I_t)>0$ and
$\mathcal I_t\subsetneq\mathcal I_{t+\Delta t}$ should be treated as properties
of the non-degenerate sector, not as axioms for every measure-dynamical
computation.  The inclusive boundary conditions are
\[
  H(\mu_t\mid\mathcal I_t)\geq 0,
  \qquad
  \mathcal I_t\subseteq\mathcal I_{t+\Delta t},
\]
with equality permitted in the classical deterministic sector.

\subsection{Exact Embedding of a Deterministic Turing Machine}

Let
\[
  M=(Q,\Gamma,\delta,q_0,q_{\mathrm{acc}},q_{\mathrm{rej}})
\]
be a deterministic Turing machine.  Its configuration space is the countable
set
\[
  \mathcal C_M
  =Q\times\Gamma_{\sqcup}^{(\mathbb Z)}\times\mathbb Z,
\]
where $\Gamma_{\sqcup}^{(\mathbb Z)}$ denotes tape configurations with finite
non-blank support and the last coordinate is the head position.  The local
transition rule $\delta$ induces a global configuration map
\[
  F_M:\mathcal C_M\longrightarrow\mathcal C_M.
\]
We totalize $F_M$ by setting $F_M(c)=c$ on halting configurations.

Equip $\mathcal C_M$ with the discrete $\sigma$-algebra and define the static
Dirac-valued Markov kernel
\[
  K_M(c,A):=\mathbf 1_A(F_M(c)).
\]
Equivalently, let the transfer operator on probability measures be the
push-forward
\[
  \mathsf P_M\mu:=(F_M)_{\#}\mu,
  \qquad
  ((F_M)_{\#}\mu)(A)=\mu(F_M^{-1}(A)).
\]

\begin{theorem}[Degenerate embedding theorem]
 For every deterministic Turing machine $M$ and every input word $w$, there
 exists a measure-dynamical computation $\mathfrak E_{M,w}$ such that, for all
 $k\in\mathbb N$,
 \[
   \mu_k=\delta_{F_M^k(c_0(w))},
 \]
 where $c_0(w)$ is the initial configuration of $M$ on $w$.  Consequently:
 \begin{enumerate}
   \item the support at each discrete time is a singleton;
   \item the Shannon entropy of the state law is zero;
   \item the transition kernel is static, deterministic, and Dirac-valued;
   \item there is no measure diffusion or endogenous update of the transition
         law; and
   \item $M$ accepts (respectively rejects) $w$ after $T$ steps if and only if
         $\mu_T$ is supported on an accepting (respectively rejecting)
         configuration.
 \end{enumerate}
 The embedding preserves discrete running time exactly: one application of
 $\mathsf P_M$ is one transition of $M$.
\end{theorem}

\begin{proof}
 Set $\mu_0=\delta_{c_0(w)}$ and recursively define
 $\mu_{k+1}=\mathsf P_M\mu_k$.  For any configuration $c$,
 \[
   (F_M)_{\#}\delta_c=\delta_{F_M(c)}.
 \]
 Hence induction on $k$ gives
 \[
   \mu_k=\delta_{F_M^k(c_0(w))}.
 \]
 A Dirac probability measure has singleton support and zero Shannon entropy.
 Since $K_M(c,\cdot)=\delta_{F_M(c)}$ for every $c$, the kernel contains no
 stochastic dispersion; since $K_M$ is independent of $k$, the transition law
 does not learn or evolve.  Finally, the accepting and rejecting claims follow
 directly from the definition of the corresponding configuration sets.  The
 recurrence applies $F_M$ exactly once per unit of discrete time, so the time
 overhead is one.
\end{proof}

This theorem gives a precise meaning to ``degenerate.''  The measures
$\delta_c$ are extreme points of the convex space $\mathcal P(\mathcal C_M)$,
and the kernels $K_M(c,\cdot)$ take values only in those extreme points.  Thus
the deterministic machine lies on the zero-entropy, zero-diffusion boundary of
the larger convex space of measure-valued transitions.

\subsection{Exact Continuous-Time Realization}

The preceding embedding is exact and requires no unjustified limit
$\Delta t\to0$.  If a continuous-time representation is desired, introduce an
internal unit clock and form the suspension state space
\[
  \widehat{\mathcal C}_M
  :=(\mathcal C_M\times[0,1])/\!\sim,
  \qquad
  (c,1)\sim(F_M(c),0).
\]
For $t\geq0$, define a semiflow by
\[
  \Phi_t([c,s])
  :=[F_M^{m}(c),s+t-m],
  \qquad m=\lfloor s+t\rfloor.
\]
The endpoint identification makes the expression well-defined and
$\Phi_{t+u}=\Phi_t\circ\Phi_u$.  Starting from
$\widehat\mu_0=\delta_{[c_0(w),0]}$, let
\[
  \widehat\mu_t=(\Phi_t)_{\#}\widehat\mu_0.
\]
Then
\[
  \widehat\mu_k=\delta_{[F_M^k(c_0(w)),0]}
  \qquad(k\in\mathbb N).
\]
Thus the continuous-time orbit reproduces the machine exactly at integer
times, and its first hitting time of a halting cross-section equals the number
of Turing steps.  At each instant the measure support is still a single point;
the union of those points over time traces a one-dimensional hybrid orbit.

The suspension is a continuous-time hybrid realization, rather than a smooth
PDE limit.  Stronger smooth realizations are known: robust analytic maps and
analytic ODE flows can simulate arbitrary Turing machines in real time.  Such
results may be cited if a smooth Euclidean realization is required, but a
complexity claim must also specify low-complexity encoders and decoders so that
the computation is not hidden in the representation map.

\begin{remark}[Nondeterministic machines]
 For a nondeterministic Turing machine, the global transition is a finite-valued
 relation $R_M(c)\subseteq\mathcal C_M$.  One may choose any kernel satisfying
 \[
   K_M(c,\{c'\})>0
   \quad\Longleftrightarrow\quad
   c'\in R_M(c).
 \]
 Then $\operatorname{supp}(\mu_k)$ is exactly the set of configurations
 reachable in $k$ nondeterministic steps, and acceptance is equivalent to the
 support intersecting the accepting set.  The weights are auxiliary: standard
 nondeterminism is existential and is encoded by support, not by probability.
 This is an atomic branching process, not automatically a high-dimensional
 continuous diffusion.
\end{remark}

\begin{remark}[Scope of the theorem]
 The theorem proves an inclusion of mathematical representations:
 deterministic Turing computations form a degenerate subclass of the proposed
 measure-dynamical framework.  It does not prove that the framework computes a
 non-Turing-computable function, that it decides an undecidable language, or
 that it solves an NP-complete problem in polynomial time.  Any claim of
 computational power strictly beyond Turing machines would require a separate
 effective model, finite-precision resource bounds, and computable encoders and
 decoders; otherwise non-computability may simply be hidden in an oracle or an
 infinitely precise real parameter.
\end{remark}

% Suggested bibliography entries:
%
% Daniel S. Graca, Manuel L. Campagnolo, and Jorge Buescu.
% Robust Simulations of Turing Machines with Analytic Maps and Flows.
% LNCS 3526, 169--179, 2005. DOI: 10.1007/11494645_21.
%
% Daniel S. Graca and Ning Zhong.
% Analytic One-Dimensional Maps and Two-Dimensional Ordinary Differential
% Equations Can Robustly Simulate Turing Machines.
% Computability 12(1), 2023. arXiv:2109.15073.
%
% Jordan Cotler and Semon Rezchikov.
% Computational Dynamical Systems. arXiv:2409.12179, 2024.
